\section{Q2. Can transfer be recovered without retraining? Output-space composition}
\label{sec:actIII}

Act I left us with an uncomfortable asymmetry. Supervised fine-tuning (SFT) buys a large,
uniform gain on the sources a guard trained on ($\RepDelta$ macro-AP,
[$\RepCILower,\RepCIUpper$]) while, on average, giving back a little transfer to sources it
never saw ($\TransferDelta$ [$\TransferCILower,\TransferCIUpper$]), and for the strongest
base the transfer cost is steep. The instinct is to tune \emph{harder} --- more data, more
steps, a different training recipe. Act~II asks the opposite question:
\emph{instead of overwriting the base's judgment, can we keep it in the decision?} The base
checkpoint, before any guard fine-tune, is often a surprisingly good \emph{transfer} scorer
(in \Cref{sec:actI} three of four bases transfer better than their own SFT guard). If SFT's
loss is that it has forgotten the base's broad, non-specialized view, then the cheapest repair
is not to retrain but to \emph{put the base back in the room} at decision time and let it vote.

\begin{background}{Why averaging two imperfect scorers can beat either}
Two guards that make \emph{different} mistakes carry complementary information. If the SFT guard
is confident-but-wrong on an off-source prompt while the base is mildly-right (or vice versa),
averaging their scores cancels part of each one's idiosyncratic error and keeps the signal they
agree on --- the same variance-reduction intuition behind any ensemble. The catch is that you
can only average two scores if they live on the \emph{same scale}; a raw margin from one model
is not comparable to a raw margin from another. That is what the calibration step below is for.
Composition is not a smarter model; it is a way to \emph{not throw away} a scorer you already
have.
\end{background}

\subsection{The fixed composition operator}
\label{sec:actIII-operator}

The rule is deliberately the simplest thing that could work. For a single input $x$ we run
\emph{both} the base and one SFT adapter, map each model's raw score to a probability with its
own calibrator, and report the fixed, equal-weight average:
\begin{equation}
s_{\mathrm{comp}}(x) \;=\; \tfrac12\, C_b\!\big(s_b(x)\big) \;+\; \tfrac12\, C_{a}\!\big(s_{a}(x)\big),
\label{eq:comp}
\end{equation}
where $s_b(x)$ and $s_a(x)$ are the base and adapter single-token margins
$z_{\text{unsafe}}-z_{\text{safe}}$ (the same head used everywhere in this report), and
$C_b,C_a$ are per-model calibrators. Three design choices in \Cref{eq:comp} are load-bearing,
and each is fixed \emph{before} looking at any transfer result.

\paragraph{Calibration makes the two scores comparable.}
A raw margin is not a probability, and two different models emit margins on two different,
incomparable scales --- a $+2.0$ from the base and a $+2.0$ from the adapter need not mean the
same confidence. Each calibrator $C(\cdot)$ is a monotone map from raw margin to a probability
in $[0,1]$, fit \emph{only on a development split} (never on the rows we score), so that after
calibration a given output value means the same ``probability this is unsafe'' for both models.
Only then does the average in \Cref{eq:comp} combine like with like rather than letting whichever
model happens to produce larger raw numbers dominate. Concretely (illustration, not a result):
if on some prompt the base calibrates to $0.30$ and the adapter to $0.90$, the composed score is
$0.60$ --- the adapter's alarm is heard but not obeyed outright.

\begin{background}{What a calibrator is, and why ``dev-split only'' matters}
A calibrator is a tiny one-input function (e.g.\ a fitted logistic or isotonic curve) that
answers ``given this raw margin, what fraction of prompts with that margin were actually
unsafe?'' It reshapes the score without reordering it --- so it changes \emph{thresholds and
comparability}, not ranking (AP is unchanged by a monotone calibrator applied to one model). We
fit it on a held-out development split so no information from the evaluation rows leaks into the
operator; this is what keeps the equal-weight average from being a disguised fit to the test set.
\end{background}

\paragraph{Equal weights, fixed in advance.}
The $\tfrac12,\tfrac12$ weights are not tuned. Choosing the mixing weight by maximizing transfer
on the very rows we then report would be circular --- it would let the operator quietly memorize
the answer. Fixing the weights a~priori means the composition has \emph{no free parameters set on
the evaluation data}; any transfer it recovers is a property of averaging base and adapter, not of
a search. (A convex-weight variant that does tune the mixing weight, at $\alpha=\PilotConvexWeight$,
was visible during development and is therefore reported only as a non-promotable ablation ---
see \Cref{sec:actIII-ablations}.)

\paragraph{Two inference passes.}
Because \Cref{eq:comp} needs both $s_b(x)$ and $s_a(x)$, composition runs \emph{two} forward
passes per input --- the base and the adapter --- roughly doubling inference cost relative to a
single guard. Nothing is retrained; there is no new checkpoint to store beyond the base and its
adapter (which a LoRA deployment already keeps). The price of skipping retraining is paid at
serving time, not training time.

\subsection{Output-space composition vs.\ weight-space merging (WiSE-FT, model soups)}
\label{sec:actIII-weightspace}

\Cref{eq:comp} combines models in \emph{output space}: it averages what they \emph{say}. The
better-known alternatives combine models in \emph{weight space}: WiSE-FT and model soups
\citep{wortsman2022wiseft,wortsman2022modelsoups} build a single merged network by interpolating
parameters, $\theta_{\mathrm{merge}}=(1-\alpha)\,\theta_{b}+\alpha\,\theta_{a}$, and then run
\emph{one} forward pass through that merged model. The two families trade off opposite costs.

\begin{itemize}[leftmargin=1.4em,itemsep=1pt]
  \item \textbf{Weight-space (WiSE-FT / soups).} One forward pass, no extra serving cost, one
  artifact to ship. But it requires the two models to live in the \emph{same, interpolable}
  weight space --- identical architecture and aligned parameters --- and it does no per-model
  calibration; you interpolate weights and hope the merged head is still well-scaled.
  \item \textbf{Output-space (ours, \Cref{eq:comp}).} Needs only \emph{comparable output scores},
  not interpolable weights, so in principle it composes models that could never be merged in
  weight space, and it calibrates each model before combining. The cost is the second inference
  pass, and the requirement that both scores be put on a common probability scale first.
\end{itemize}

We test only the output-space operator here. Because the two approaches optimize different
constraints, output-space recovery does \emph{not} predict weight-space recovery, and we do not
claim it does; a direct WiSE-FT rescoring control is a stated gap (\Cref{sec:actIII-ablations}).

\subsection{Results: composition recovers transfer at a small represented cost}
\label{sec:actIII-results}

\Cref{tab:pilot-summary} reports the fixed-panel macro-AP of each operator in both regimes, plus
the worst-of-the-two-regimes column $\min(\text{both})$. The pattern is clean. The base is the
better of the two \emph{single} guards on transfer ($\PilotBaseTransfer$) but a poor represented one
($\PilotBaseRepresented$). SFT inverts that: it is the best \emph{represented} scorer
($\PilotSFTRepresented$) but the worst transfer scorer ($\PilotSFTTransfer$) --- the Act~I
specialization, restated. Composition sits between the two on represented ranking
($\PilotCompositionRepresented$) and \emph{above the base} on transfer
($\PilotCompositionTransfer$). Read down the $\min(\text{both})$ column, composition is the most
\emph{balanced} promotable operator: its worst regime ($0.883$) beats both the base's worst
($0.658$) and SFT's worst ($0.807$). In other words, if you must commit to one scorer without
knowing whether the next prompt is on-source or off-source, the composition is the best-balanced single
choice on this panel by worst-regime AP --- a \emph{ranking} statement, not a deployability one (its
operating point still misses the illustrated FPR target; see below).

% Generated by code/build_pilot_artifacts.py; do not edit by hand.
\begin{table}[H]
\centering
\caption{Retrospective clean-v2 fixed-panel macro-AP. The calibrated average is the fixed primary operator. Logit averaging is an exposed ablation and cannot be promoted after observing transfer results.}
\label{tab:pilot-summary}
\small
\begin{tabular}{lrrr}
\toprule
Guard & Represented & Transfer & $\min$(both) \\
\midrule
Unadapted base & 0.658 & 0.866 & 0.658 \\
SFT adapter & 0.982 & 0.807 & 0.807 \\
Base+SFT calibrated average & 0.962 & 0.883 & 0.883 \\
Base+SFT logit average (ablation) & 0.943 & 0.891 & 0.891 \\
\bottomrule
\end{tabular}
\end{table}


Against SFT --- the relevant baseline, since composition is a repair \emph{for} an SFT guard ---
composition changes represented-source macro-AP by $\PilotRepDeltaSFT$
[$\PilotRepDeltaSFTLow,\PilotRepDeltaSFTHigh$] and transfer by $\PilotTransferDeltaSFT$
[$\PilotTransferDeltaSFTLow,\PilotTransferDeltaSFTHigh$]. That is the headline trade: it gives
back under two points of represented ranking to buy back roughly eight points of transfer.
Against the untuned base, the aggregate transfer gain is smaller and its interval is closer to
zero, $\PilotTransferDeltaBase$ [$\PilotTransferDeltaBaseLow,\PilotTransferDeltaBaseHigh$]:
composition edges past the base \emph{on average}, but only for $\PilotNTransferPositiveVsBase$ of
the four checkpoints individually --- and that interval resamples rows and seeds only, so the edge sits
inside the $0.015$--$0.029$ reproduction envelope measured in \Cref{sec:actI-klsft} and the
vs.-base direction is \emph{unresolved}. All intervals here are descriptive paired percentile-bootstrap
ranges under the retrospective, estimation-only regime (\PilotEvidenceStatus{}), conditional on
this fixed panel --- not significance tests and not population claims.

\subsubsection{Per-checkpoint recovery}
\label{sec:actIII-percheckpoint}

The aggregate hides a more instructive per-checkpoint story, in \Cref{tab:pilot-models} and
\Cref{fig:act3-bars}.

% Generated by code/build_pilot_artifacts.py; do not edit by hand.
\begin{table}[H]
\centering
\caption{Transfer macro-AP by checkpoint. Deltas are observed contrasts with descriptive paired-bootstrap percentile intervals. Recovery relative to SFT is positive for every checkpoint, but comparison with base is heterogeneous.}
\label{tab:pilot-models}
\small
\resizebox{\linewidth}{!}{%
\begin{tabular}{lrrrrr}
\toprule
Checkpoint & Base & SFT & Base+SFT & $\Delta$ vs. SFT [95\% CI] & $\Delta$ vs. base [95\% CI] \\
\midrule
SmolLM2-1.7B & 0.790 & 0.830 & 0.857 & +0.027 [+0.013, +0.043] & +0.067 [+0.038, +0.095] \\
Qwen2.5-1.5B & 0.819 & 0.780 & 0.855 & +0.075 [+0.047, +0.103] & +0.036 [+0.008, +0.066] \\
SmolLM3-3B & 0.910 & 0.823 & 0.907 & +0.084 [+0.063, +0.104] & -0.003 [-0.012, +0.005] \\
Qwen3-4B & 0.944 & 0.794 & 0.914 & +0.120 [+0.082, +0.160] & -0.030 [-0.043, -0.018] \\
\bottomrule
\end{tabular}%
}
\end{table}


\begin{figure}[t]\centering
\includegraphics[width=0.82\linewidth]{fig_act3_composition.pdf}
\caption{Per-checkpoint transfer macro-AP for base, SFT, and the base+SFT composition. Composition (green)
exceeds SFT (orange) for every checkpoint; relative to the base it lands above for the two weaker bases,
approximately at base for SmolLM3-3B ($0.907$ vs.\ $0.910$), and below base for the strongest base
(Qwen3-4B) --- transfer \emph{recovery}, not Pareto dominance.}
\label{fig:act3-bars}
\end{figure}

\begin{itemize}[leftmargin=1.4em,itemsep=2pt]
  \item \textbf{SmolLM2-1.7B} ($0.790\!\to\!0.830\!\to\!0.857$): the friendly case. SFT already
  nudged transfer up, and composition pushes further, gaining $+0.027$ over SFT and $+0.067$ over
  base --- both regimes end up ahead.
  \item \textbf{Qwen2.5-1.5B} ($0.819\!\to\!0.780\!\to\!0.855$): the clearest \emph{repair}. Here
  SFT actively \emph{hurt} transfer (an Act~I specialization loss); composition not only reverses
  it ($+0.075$ vs.\ SFT) but climbs back above the base ($+0.036$).
  \item \textbf{SmolLM3-3B} ($0.910\!\to\!0.823\!\to\!0.907$): composition almost exactly
  \emph{undoes} the SFT damage ($+0.084$ vs.\ SFT), landing essentially back at the base
  ($-0.003$, interval $[-0.012,+0.005]$ straddling zero) --- recovery to the base, not past it.
  \item \textbf{Qwen3-4B} ($0.944\!\to\!0.794\!\to\!0.914$): the recurring character. It was the
  worst specializer in Act~I, so composition helps it \emph{most relative to SFT} ($+0.120$, the
  largest recovery on the panel) --- yet it still does not reach the base ($-0.030$, interval
  $[-0.043,-0.018]$, entirely below zero). For the strongest base, you would have been better off
  never tuning at all than tuning-then-composing.
\end{itemize}

\begin{takeaway}{Recovery, not Pareto dominance}
Composition beats SFT on transfer for \emph{all four} checkpoints (every ``$\Delta$ vs.\ SFT''
interval sits above zero), so as a remedy for an already-specialized guard it is reliable on this
panel. But it does not dominate the base: vs.\ base the four deltas are heterogeneous
($+0.067,+0.036,-0.003,-0.030$) --- helping two, neutral on one, and \emph{hurting} the strongest
base (Qwen3-4B). The honest reading is therefore ``composition recovers much of the transfer SFT
gave up,'' not ``composition is free improvement over doing nothing.'' If transfer is what you
care about and you have not yet tuned, the base alone can still be the better scorer.
\end{takeaway}

\subsubsection{Why composition helps at all --- and least where the base is strongest}
\label{sec:actIII-why}

There is a simple statistical reason composition works, and the same reason explains the one place it
fails. Averaging two scorers is an \emph{ensemble}, and an ensemble beats its members when they are each
individually good and make \emph{different} mistakes. A clean diagnostic is the \emph{midpoint}: if
composition merely interpolated between the base and the SFT guard, it would land at their average AP. It
does not --- it lands \emph{above} that midpoint by a consistently positive margin on every checkpoint
($+0.047,+0.055,+0.040,+0.045$; read off \Cref{tab:pilot-models}). Because AP is nonlinear, this
above-midpoint margin is a heuristic \emph{signature} of diversity, not a formal decomposition; a direct
measurement bears it out --- on transfer the base's per-row errors correlate only \EnsErrCorrBaseSFT{} with
its own fine-tune's, versus \EnsErrCorrSFTSFT{} between two fine-tune seeds (\Cref{app:ensembling}). The
base and its own fine-tune rank the hard transfer cases differently, and averaging cancels part of each
one's error --- which is why composition is not mere interpolation.

The same lens explains why composition helps \emph{least} where the base is strongest. The gain from
averaging shrinks as the two scorers' errors grow more \emph{correlated}. A LoRA adapter is a low-rank edit
\emph{anchored} to its base, so a strong base's fine-tune stays close to it: their errors correlate more,
the diversity term shrinks, and the average can no longer clear the (already high) base. That is exactly the
Qwen3-4B story --- strongest base, most base-correlated adapter, composition landing below base. Acts~I and
II are then one mechanism seen twice: the strongest base specializes \emph{most} (Act~I) and is helped
\emph{least} by composition (Act~II), both because its fine-tune moves the least far from it.

\subsubsection{The equal-cost control: it is the base that helps, not a second scorer}
\label{sec:actIII-sftsft}

The diversity account makes a falsifiable prediction: if the recovery came from generic two-model
ensembling, then averaging \emph{two SFT adapters} --- same two-pass inference cost, but no base --- should
recover about as much; if instead it is the base's less-specialized view that matters, the SFT$+$SFT
average should recover far less. This control needs \emph{no new training}: we already hold five scored SFT
seeds per checkpoint, so composing two of them is a pure recompute of the committed calibrated per-row
scores. \Cref{tab:sftsft} runs it. base$+$SFT beats SFT$+$SFT on \emph{every} checkpoint, and the gap is
largest exactly where the base is strongest ($+0.066$ for SmolLM3-3B, $+0.102$ for Qwen3-4B, vs.\ $+0.013$
for the already-friendly SmolLM2-1.7B). Two independently seeded adapters are near-copies of one another, so
their average barely improves on a single SFT guard ($0.79$--$0.84$, close to SFT's own transfer); the base
contributes a genuinely different ranking of the hard off-source cases, and that is what the composition is
cashing in. The recovery is therefore attributable to \emph{keeping the base}, not to running a second model.

\input{generated/tab_sftsft_gen}

\subsection{Ranking is not calibration: the operating-point gap}
\label{sec:actIII-operating}

Everything above is about \emph{ranking} (AP), which needs no threshold. Deployment needs a
threshold, and here composition's win is only partial. \Cref{tab:pilot-operating} sets each guard's
threshold for a $\PilotTargetFPRPct$\% false-positive target on calibration negatives and reports
the \emph{realized} transfer rates. Composition improves recall (macro-TPR $0.517\!\to\!0.639$
across base$\to$comp, best of the three) and its realized transfer false-alarm rate,
$\PilotCompositionTransferFPRPct$\%, is much better than SFT's $\PilotSFTTransferFPRPct$\% --- but
it still overshoots the $\PilotTargetFPRPct$\% target and remains \emph{above the base's}
$\PilotBaseTransferFPRPct$\%. The pooled-FPR column tells the same story ($0.043/0.170/0.091$ for
base/SFT/comp). Recovering \emph{rank} did not hand us a threshold that \emph{transfers}: the
cutoff chosen on calibration data lands in the wrong place off-source, because the score
\emph{distribution} shifts between the calibration and transfer data even when the \emph{ordering}
recovers. So treat an AP recovery as a reason to \emph{re-calibrate on the target regime}, never as
evidence that the old threshold is safe to reuse.

% Generated by code/build_pilot_artifacts.py; do not edit by hand.
\begin{table}[H]
\centering
\caption{Realized transfer operating points after choosing thresholds for a 5\% FPR target on calibration negatives. Rank recovery does not imply calibration transfer.}
\label{tab:pilot-operating}
\small
\begin{tabular}{lrrr}
\toprule
Guard & Macro TPR & Macro FPR & Pooled FPR \\
\midrule
Unadapted base & 0.517 & 0.081 & 0.043 \\
SFT adapter & 0.581 & 0.155 & 0.170 \\
Base+SFT calibrated average & 0.639 & 0.114 & 0.091 \\
\bottomrule
\end{tabular}
\end{table}


\subsection{Ablations, and what Act II does \emph{not} establish}
\label{sec:actIII-ablations}

\paragraph{The logit-average ablation (non-promotable).}
The last row of \Cref{tab:pilot-summary} reports averaging the two models' \emph{raw margins}
before calibration (a logit-space average) rather than the calibrated probabilities of
\Cref{eq:comp}. It edges the calibrated operator on transfer ($\PilotLogitTransfer$ vs.\
$\PilotCompositionTransfer$). We nonetheless \emph{do not promote it}: both it and the convex-weight
variant ($\alpha=\PilotConvexWeight$) were visible while we were developing the method, so choosing
the operator after seeing that it wins on transfer would be exactly the retrospective cherry-pick
this report is built to avoid. The calibrated equal-weight average was fixed in advance and is the
only primary operator; the alternatives are reported for transparency, as ablations, with no claim
attached.

\paragraph{Stated gaps.}
This is a pilot, and one control that would sharpen the mechanism is still missing.
\emph{A real WiSE-FT rescoring control.} We contrast against weight-space merging conceptually
(\Cref{sec:actIII-weightspace}) but do not actually rescore a WiSE-FT / soup checkpoint on this panel, so
we cannot say which family recovers more transfer here. (The other control we flagged in earlier drafts ---
an equal-cost SFT$+$SFT average --- is now \emph{run} rather than pending, in \Cref{sec:actIII-sftsft}: it
confirms the recovery is the base's doing, not generic ensembling.)

\paragraph{And what it does not license.}
No Pareto-dominance claim (composition can leave the strongest base worse off); no causal or
mechanistic claim (the ensemble intuition is a motivation, not a proof); and no population or
significance claim (all numbers are descriptive contrasts on a fixed four-checkpoint panel, with a
manifest inspected during development). Finally, we compose only the SFT adapter here; whether the same
operator recovers transfer for guards tuned with \emph{other objectives} (e.g.\ preference optimization)
is left open.

\edb{Output-space base$+$SFT averaging recovers transfer relative to SFT ($\PilotTransferDeltaSFT$) and
beats an equal-cost SFT$+$SFT ensemble (\Cref{tab:sftsft}), so the gain is the base's, not generic
ensembling; it has the highest worst-regime AP of the \emph{promotable} operators (the dev-visible
  logit-average ablation is nominally better still, $0.891$; \Cref{sec:actIII-ablations}).}
{If you have already SFT'd and transfer regressed, compose base$+$adapter to recover it without
retraining --- then recalibrate the threshold on the target regime.}
{Ranking recovery $\neq$ calibration transfer: the composition operating point still misses the
illustrated $5\%$ FPR target ($11.4\%$). Descriptive, fixed-panel; the ensemble intuition is a
motivation, not a proof.}
